% Responsable: P1  (versión completa en extendido/10_conclusiones.tex)
\section{Conclusiones}
\label{sec:conclusiones}

\begin{itemize}
  \item \textbf{Una CFG describe bien un fragmento del español, y solo ese fragmento.} Dentro de
    él, el parser es exacto en decisión, tipo de error y análisis; fuera de él, el 20\,\% del
    español correcto de nuestro conjunto queda rechazado sin análisis.
  \item \textbf{La concordancia es lo que aporta la CFG con rasgos.} Es el único sistema que
    rechaza las 13 oraciones con error de concordancia, dice por qué y no da falsas alarmas.
    Verificarla aparte, como en la fase semántica de un compilador, fue lo que la hizo manejable.
  \item \textbf{La ambigüedad es el límite común.} La gramática no es LL(1) por la frase
    preposicional; el parser acierta la unión en el 59\,\% de los casos, spaCy y el LLM en más,
    pero ninguno reconoce las oraciones ambiguas.
  \item \textbf{Los modelos estadísticos ganan en cobertura y pierden en garantías:} etiquetan
    casi todo, pero no deciden la gramaticalidad o la confunden con frecuencia.
  \item \textbf{Un LLM puede seguir reglas formales, pero no se puede confiar en que lo haga:}
    decide como las reglas (0.97--0.99), pero a veces lo justifica con árboles inválidos.
\end{itemize}

\paragraph{Limitaciones y trabajo futuro.} El conjunto tiene 73 oraciones etiquetadas por el
mismo equipo que escribió la gramática. El parser no retrocede, así que pierde algunas
oraciones con palabras ambiguas; un parser de gráfico resolvería eso y mostraría todas las
lecturas de las oraciones ambiguas. La gramática se puede ampliar con sujeto omitido y
coordinación, y el experimento con el LLM, repetirse con modelos más grandes.
